\section{Komputer Kuantum}

\subsection{Sistem Qubit dan Geometri Bola Bloch}
Komputer kuantum merepresentasikan unit informasi terkecil melalui konsep \textit{qubit} yang memperluas paradigma bit klasik menjadi sistem dua-keadaan mekanika kuantum yang jauh lebih kompleks. Secara matematis, sebuah \textit{qubit} didefinisikan sebagai vektor satuan yang beroperasi di dalam ruang Hilbert kompleks dua-dimensi $\mathbb{C}^2$, yang memungkinkan representasi informasi dalam skala probabilitas yang kontinu. Keadaan umum dari sebuah \textit{qubit} $|\psi\rangle$ dinyatakan sebagai superposisi linear dari basis komputasi $|0\rangle$ dan $|1\rangle$, yang mencerminkan kemungkinan eksistensi simultan dari status biner yang berlawanan. Parameterisasi vektor keadaan ini dilakukan menggunakan sudut polar $\theta$ dan fase azimuthal $\phi$ guna memetakan status \textit{qubit} ke dalam permukaan geometris yang intuitif bagi para peneliti.

Visualisasi geometris terhadap keadaan \textit{qubit} dilakukan melalui representasi bola Bloch, di mana setiap titik pada permukaan bola satuan 3D tersebut berkorespondensi dengan satu vektor keadaan yang unik. Vektor keadaan yang dipetakan ke permukaan bola Bloch memberikan kemudahan dalam memahami proses rotasi parameter yang terjadi selama operasional algoritma kuantum berlangsung. Penggunaan geometri bola Bloch ini sangat relevan dalam domain finansial karena membantu memetakan keputusan biner aset—seperti pilihan untuk membeli atau menjual—ke dalam ruang parameter kontinu yang dapat dioptimasi secara presisi. Dengan melakukan rotasi pada vektor keadaan, sistem dapat mengeksplorasi berbagai tingkat kepercayaan atau probabilitas strategi investasi sebelum akhirnya dilakukan pengukuran untuk mendapatkan hasil akhir.

\subsection{Prinsip Superposisi dan Amplitudo Probabilitas}
Prinsip superposisi merupakan karakteristik fundamental mekanika kuantum yang memberikan kemampuan bagi \textit{qubit} untuk berada dalam banyak keadaan basis secara simultan sebelum proses pengukuran dilakukan. Amplitudo probabilitas, yang direpresentasikan oleh koefisien kompleks $\alpha$ dan $\beta$, menyimpan informasi vital mengenai fase koherensi serta peluang kemunculan dari masing-masing keadaan basis saat sistem diobservasi. Sifat superposisi ini melahirkan fenomena \textit{quantum parallelism}, di mana sebuah sistem dengan $n$-\textit{qubit} mampu merepresentasikan $2^n$ konfigurasi solusi secara simultan dalam satu vektor keadaan tunggal. Dalam optimasi portofolio, paralelisme ini memungkinkan komputer kuantum untuk mengevaluasi seluruh kombinasi aset yang mungkin dalam satu waktu, memberikan keunggulan komputasi yang signifikan dibandingkan metode klasik.

Transisi dari keadaan superposisi menuju keadaan deterministik terjadi saat dilakukan pengukuran, yang mengakibatkan runtuhnya fungsi gelombang (\textit{wavefunction collapse}) menjadi satu keadaan pasti berdasarkan Aturan Born. Peluang untuk menemukan sistem pada hasil pengukuran tertentu ditentukan secara eksak oleh nilai kuadrat modulus dari amplitudo probabilitas yang terkait dengan keadaan basis tersebut. Sifat probabilistik intrinsik ini bukan merupakan batasan, melainkan sebuah fitur yang memungkinkan eksplorasi ruang solusi secara non-deterministik pada lanskap energi finansial yang sangat kompleks. Dengan memanfaatkan kolaps fungsi gelombang, algoritma optimasi dapat menarik solusi yang paling efisien dari distribusi probabilitas yang telah dibentuk melalui proses evolusi uniter pada sirkuit kuantum.

\subsection{Keterbelitan Kuantum (\textit{Entanglement})}
Keterbelitan kuantum (\textit{entanglement}) merupakan fenomena non-klasik di mana keadaan satu \textit{qubit} tidak dapat didefinisikan secara independen dari keadaan \textit{qubit} lainnya, meskipun keduanya dipisahkan oleh jarak fisik. Dalam konteks pemodelan portofolio finansial, keterbelitan mewakili hubungan ketergantungan non-lokal antar variabel aset yang secara matematis melampaui kemampuan matriks korelasi linear tradisional. Melalui fenomena ini, komputer kuantum dapat menangkap struktur kovariansi aset yang sangat kompleks, di mana keputusan pada satu aset secara instan mempengaruhi probabilitas keputusan pada aset lainnya yang saling terikat. Representasi keadaan terbelit maksimal, yang dikenal sebagai \textit{Bell States}, menjadi fondasi utama bagi transmisi informasi kuantum yang efisien serta pembangunan sirkuit yang mampu memodelkan sistem keuangan global secara utuh.

\subsection{Gerbang Logika dan Evolusi Uniter}
Manipulasi informasi di dalam komputer kuantum dilakukan melalui gerbang logika kuantum yang merupakan operator unitari untuk mentransformasi vektor keadaan di dalam ruang Hilbert. Setiap operator gerbang harus memenuhi syarat unitaritas yang ketat, yaitu $U^\dagger U = I$, guna menjaga integritas informasi melalui sifat reversibilitas serta menjamin konservasi total probabilitas sistem. Operasi gerbang kuantum mengikuti prinsip evolusi yang diatur oleh persamaan Schrödinger, yang mewakili evolusi waktu dari sistem kuantum yang terkontrol secara presisi melalui interaksi fisik. Karakteristik uniter ini memastikan bahwa setiap proses komputasi yang dijalankan pada sirkuit kuantum tidak akan kehilangan informasi dan tetap berada pada permukaan bola Bloch yang merepresentasikan keadaan fisik yang sah.

Implementasi fisik dari gerbang kuantum pada perangkat keras dilakukan melalui pemberian pulsa elektromagnetik yang dikontrol dengan presisi tinggi untuk memutar vektor keadaan pada sumbu-sumbu bola Bloch. Klasifikasi gerbang kuantum terdiri dari gerbang satu-\textit{qubit} yang berfungsi untuk melakukan rotasi fase dan menciptakan superposisi, serta gerbang multi-\textit{qubit} yang berperan dalam membangun keterbelitan antar partikel. Kombinasi sistematis dari berbagai jenis gerbang universal ini membentuk arsitektur sirkuit kuantum yang secara kolektif merepresentasikan fungsional biaya finansial dalam bentuk Hamiltonian energi. Dengan menyusun sirkuit dari gerbang-gerbang ini, peneliti dapat merancang algoritma yang mampu menavigasi kompleksitas pasar finansial melalui manipulasi keadaan kuantum yang terukur dan efisien secara operasional.
